\chapter{יישומים והערכת מודלים}

פרק זה עוסק ביישומים מעשיים של מודלים מבוססי \textenglish{Transformer}
בתחומים שונים, לצד שיטות הערכה ומדדי ביצוע מקובלים בתחום.

\section{כוונון עדין}

כוונון עדין (\textenglish{Fine-Tuning}) הוא תהליך שבו מודל מאומן-מקדים
(\textenglish{Pre-trained Model}) מותאם למשימה ספציפית. גישת \textenglish{Transfer Learning}
מאפשרת ניצול ייצוגים שנלמדו על קורפוסים גדולים למשימות עם נתונים מוגבלים.

\subsection{\textenglish{BERT} לסיווג טקסט}

\textenglish{BERT} (\textenglish{Bidirectional Encoder Representations from Transformers})
הוא מודל מקודד דו-כיווני שאומן על משימות \textenglish{Masked Language Modeling}
ו-\textenglish{Next Sentence Prediction}.

לכוונון עדין לסיווג, מוסיפים ראש סיווג (\textenglish{Classification Head})
על גבי אסימון \textenglish{[CLS]}:

\begin{equation}
    P(y \mid x) = \text{softmax}(W \cdot h_{\text{[CLS]}} + b)
\end{equation}

כאשר \(h_{\text{[CLS]}}\) הוא ייצוג האסימון הראשון בפלט \textenglish{BERT}.

\subsection{\textenglish{GPT} ליצירת טקסט}

מודלי \textenglish{GPT} (\textenglish{Generative Pre-trained Transformer})
מבוססים על מפענח (\textenglish{Decoder}) בלבד ומאומנים על מידול שפה אוטו-רגרסיבי:

\begin{equation}
    P(w_1, \ldots, w_n) = \prod_{i=1}^{n} P(w_i \mid w_1, \ldots, w_{i-1})
\end{equation}

\section{מדדי הערכה}

\subsection{דיוק ורקול}

מדדי הדיוק (\textenglish{Precision}), הרקול (\textenglish{Recall}),
ו-\textenglish{F1} הם הבסיסיים ביותר לסיווג:

\begin{equation}
    \text{Precision} = \frac{TP}{TP + FP}, \qquad
    \text{Recall}    = \frac{TP}{TP + FN}
\end{equation}

\begin{equation}
    F_1 = 2 \cdot \frac{\text{Precision} \cdot \text{Recall}}
                       {\text{Precision} + \text{Recall}}
\end{equation}

\subsection{\textenglish{BLEU} לתרגום מכונה}

מדד \textenglish{BLEU} (\textenglish{Bilingual Evaluation Understudy})
משמש להערכת תרגום מכונה. הוא מחשב חפיפה של \textenglish{n-gram}
בין תרגום מוצע לתרגום ייחוס:

\begin{equation}
    \text{BLEU} = BP \cdot \exp\!\left( \sum_{n=1}^{N} w_n \log p_n \right)
\end{equation}

כאשר \(BP\) הוא עונש קיצור (\textenglish{Brevity Penalty}) ו-\(p_n\) הוא
דיוק \textenglish{n-gram} משוקלל.

\subsection{\textenglish{Perplexity}}

מדד ה-\textenglish{Perplexity} מודד את איכות מודל שפתי:

\begin{equation}
    \text{PPL}(W) = P(w_1, w_2, \ldots, w_N)^{-\frac{1}{N}}
\end{equation}

ערך נמוך יותר של \textenglish{Perplexity} מצביע על מודל טוב יותר.

\section{דוגמת קוד: כוונון עדין של \textenglish{BERT}}

\begin{LTR}
\begin{verbatim}
from transformers import BertForSequenceClassification, BertTokenizer
from torch.utils.data import DataLoader
import torch

tokenizer = BertTokenizer.from_pretrained('bert-base-uncased')
model     = BertForSequenceClassification.from_pretrained(
                'bert-base-uncased', num_labels=2)

def encode_batch(texts, labels):
    enc = tokenizer(texts, padding=True, truncation=True,
                    return_tensors='pt', max_length=512)
    enc['labels'] = torch.tensor(labels)
    return enc

optimizer = torch.optim.AdamW(model.parameters(), lr=2e-5)

for epoch in range(3):
    for batch in dataloader:
        outputs = model(**batch)
        loss    = outputs.loss
        loss.backward()
        optimizer.step()
        optimizer.zero_grad()
    print(f"Epoch {epoch+1} done, loss={loss.item():.4f}")
\end{verbatim}
\end{LTR}

\section{השוואת מודלים}

\begin{english}
Modern language models are evaluated on a variety of benchmarks. The
\textenglish{GLUE} (\textenglish{General Language Understanding Evaluation}) and
\textenglish{SuperGLUE} benchmarks measure performance across diverse NLP tasks.
State-of-the-art models such as \textenglish{GPT-4} and \textenglish{LLaMA-3}
achieve near-human performance on many of these tasks.
\end{english}

\section{אתגרים ומגמות עתידיות}

למרות ההצלחות המרשימות, מודלי שפה גדולים עדיין מתמודדים עם אתגרים:

\begin{itemize}
    \item הזיות (\textenglish{Hallucinations}) — יצירת עובדות שגויות
    \item הטיות (\textenglish{Biases}) — שיקוף הטיות קיימות בנתוני האימון
    \item עלויות חישוב (\textenglish{Computational Cost}) — דרישות חומרה גבוהות
    \item פרשנות (\textenglish{Interpretability}) — קושי בהבנת תהליך ההחלטה
\end{itemize}

מגמות מחקר עתידיות כוללות \textenglish{Retrieval-Augmented Generation}
(\textenglish{RAG}), \textenglish{Chain-of-Thought Prompting}, ופיתוח
מודלים יעילים יותר מבחינת צריכת אנרגיה.

\section{סיכום}

פרק זה סיכם את היישומים המרכזיים של מודלים מבוססי \textenglish{Transformer},
שיטות כוונון עדין, ומדדי הערכה מקובלים. העבודה האקדמית שהוצגה לאורך ספר זה
מדגימה את הפוטנציאל העצום של למידה עמוקה ועיבוד שפה טבעית ב-\textenglish{AI}
המודרני.
